\section{Extended Analysis}

\subsection{Surprise as a detector of shift}

The allocation policy assumes that surprise increases when the model is evaluated outside its
training distribution. We can test the assumption directly, without any planner: for each
suite, encode held-out episodes in distribution and under shift, compute the one-step
Mahalanobis residual in latent space, and compare the distributions. The model is trained the
same way in both cases; only the evaluation stream changes. This is the reward-free analogue
of LeWM's "surprise evaluation", and it is a pre-condition for using surprise as a control
signal.

\subsection{Allocation is non-uniform}

If the controller were useless, allocated steps and refinement depth would be constant across
episodes. Instead, the statistics collected during evaluation (Figure~\ref{fig:stats}) show a
spread of allocated steps, learning-rate multipliers and refinement depths, with the mean
shifted upwards on shifted suites. The controller's behaviour is additionally logged per
replan, so the relationship between relative surprise and allocated compute can be read off
directly from \texttt{results/raw.jsonl}.

\subsection{Relation to AdaJEPA's hand-tuned adaptation}

AdaJEPA's own ablations show that for out-of-distribution settings, "stronger adaptation---via
a larger learning rate, more optimization steps, or a larger recent replay buffer---is often
beneficial when latency permits", while in distribution a single step at the training learning
rate is the practical default. That is exactly a two-point policy. DTA-JEPA can be read as
learning the \emph{shape} of that policy from a signal the agent already computes for free---the
residual of its own prediction---instead of selecting one of two points by hand per
environment family, and additionally spending recursive depth, which AdaJEPA's architecture
does not expose.

\subsection{Failure modes we observed}

\begin{itemize}[leftmargin=1.2em,itemsep=2pt]
\item \textbf{Symmetric shapes are easier than they look.} The square and I blocks have
$\pi/2$ rotational symmetry, so the orientation tolerance is met on a larger fraction of the
pose space. Success rates on those shapes should not be read as geometric generalisation.
\item \textbf{Colour shifts that remove the agent--distractor distinction are hardest.}
Rendering the static reference marker in the same red as the moving block removes the visual
cue that separates "what moves" from "what does not", and no amount of predictor adaptation
recovers it, consistent with AdaJEPA's finding that red-anchor and red-block shifts are where
their gains are smallest.
\item \textbf{Replay-buffer entropy.} Because we execute five actions per replan, the recent-5
buffer covers one replanning step of experience, and its entropy is low. Memory-based
retrieval is therefore most useful late in an episode or across episodes, which is what the
skill bank targets.
\end{itemize}
